
\subsubsection{Parameter-Efficient Fine-Tuning}

Low-Rank Adaptation (LoRA) introduced rank-decomposition matrices injected into transformer layers, achieving comparable performance to full fine-tuning with less than 1\% trainable parameters. AdapterFusion demonstrated that task-specific adapters can be combined non-destructively through learned attention, outperforming multi-task learning on 16 NLU tasks. These methods establish that adapter combination is both mathematically valid (LoRA deltas are additive) and empirically effective, but require knowing which adapters to combine for a given input.

\subsubsection{Dynamic Adapter Routing}

LoRAHub achieves cross-task generalization via gradient-free composition, optimizing adapter weights over few-shot validation loss. The approach achieves upper-bound performance exceeding in-context learning on Big-Bench Hard, but requires validation examples to compute the optimization objective.

LORAUTER routes queries using task embeddings computed from validation sets, achieving 101.2\% of oracle performance on task-aligned adapters and +5.2 points on unseen tasks. The method scales to 1500+ adapters but relies on 5+ validation examples per task to construct the routing representation.

MoELoRA treats LoRA modules as MoE experts with contrastive learning to encourage distinct specialization, achieving +4.2\% over vanilla LoRA on math reasoning. However, routing is jointly trained with adapters, precluding generalization to new tasks without retraining.

Multi-Head Adapter Routing (MHR) combines adapter subsets with learned routing weights. HMoRA extends this to hierarchical routing with auxiliary losses.

All existing routing methods either require validation data, joint training, or gradient signals at test time. None test whether instruction-adapter alignment exists intrinsically without any task-specific training signal.

\subsubsection{Instruction Tuning and Task Representations}

FLAN demonstrated that instruction-tuned models condition behavior on instruction format across diverse task families. T0 showed zero-shot generalization through unified prompting. These works suggest that instruction prefixes encode meaningful task semantics.

Sentence encoders like MiniLM and Sentence-BERT compress text into fixed-dimensional embeddings via contrastive pretraining. SimCSE demonstrated that such embeddings provide paraphrase invariance---a property tested explicitly in the robustness experiments.
